\documentclass[11pt]{article}
\usepackage{array}
\usepackage{booktabs}
\usepackage[T1]{fontenc}
\usepackage[utf8]{inputenc}
\usepackage{amsfonts,url,epsfig,float,xcolor}
%%% Document layout, margins
\usepackage{geometry} 
\geometry{letterpaper, textwidth=6.5in, textheight=9in, marginparsep=1em}
%%% Section headings
\usepackage{sectsty} 
\usepackage[normalem]{ulem}
%\setlength{\baselineskip}{40pt}
%\usepackage{pdfpages}

\sectionfont{\sffamily\bfseries\upshape\large}
\subsectionfont{\sffamily\bfseries\upshape\normalsize} 
\subsubsectionfont{\sffamily\mdseries\upshape\normalsize}
\makeatletter
\renewcommand\@seccntformat[1]{\csname the#1\endcsname.\quad}
\makeatother\renewcommand{\bibitem}{\vskip 2pt\par\hangindent\parindent\hskip-\parindent}
\newcommand{\mme}{\mathbb{E}}

\makeatletter
\def\@maketitle{%
  \begin{center}%
  \let \footnote \thanks
    {\large \@title \par}%
    {\normalsize
      \begin{tabular}[t]{c}%
        \@author
      \end{tabular}\par}%
    {\small \@date}%
  \end{center}%
}
\makeatother


\usepackage{mathtools}
\usepackage{extarrows}
\usepackage{enumerate} 
\usepackage{hyperref}
\usepackage{bbm}
\hypersetup{
	colorlinks=true,
	linkcolor=black,
	citecolor=black,
	filecolor=black,
	urlcolor=black,
}
\usepackage{amsthm}
\usepackage{amsmath}
\usepackage{amssymb}
\newtheorem{theorem}{Theorem}
\newtheorem{lemma}{Lemma}
\newtheorem{proposition}{Proposition}
\usepackage[toc,page]{appendix}
\newcommand{\note}[1]{{\color{red}#1}}
 \usepackage{comment}
 \usepackage{graphicx}
\usepackage{wrapfig}
\usepackage{bbm}
\usepackage{microtype}
 
\usepackage[ruled,vlined]{algorithm2e}
\usepackage{relsize}
\let\oldv\verbatim
\let\oldendv\endverbatim
\def\verbatim{\par\setbox0\vbox\bgroup\oldv}
\def\endverbatim{\oldendv\egroup\fboxsep0pt \noindent\colorbox[gray]{0.96}{\usebox0}\par}
%%%%%
 % \startlocaldefs

 % \endlocaldefs

\DeclareMathOperator{\SNR}{\mathrm{SNR}}
\DeclareMathOperator{\normal}{\mathrm{normal}}
\DeclareMathOperator{\MVN}{\mathrm{MVN}}
\DeclareMathOperator{\uniform}{\mathrm{uniform}}
\DeclareMathOperator{\Poisson}{\mathrm{Poisson}}
\DeclareMathOperator{\exponential}{\mathrm{exponential}}
\DeclareMathOperator{\binomial}{\mathrm{binomial}}
\DeclareMathOperator{\Cauchy}{\mathrm{Cauchy}}
\DeclareMathOperator{\logit}{\mathrm{logit}}
\DeclareMathOperator{\logistic}{\mathrm{logistic}}
\DeclareMathOperator{\lognormal}{\mathrm{lognormal}}
\DeclareMathOperator{\E}{\mathrm{E}}
\DeclareMathOperator{\Var}{\mathrm{var}}
\DeclareMathOperator{\sd}{\mathrm{sd}}
\DeclareMathOperator{\Cov}{\mathrm{cov}}
\DeclareMathOperator{\Dev}{\mathrm{Dev}}
\DeclareMathOperator{\Beta}{\mathrm{beta}}
\DeclareMathOperator{\mode}{\mathrm{mode}}
\renewcommand{\Pr}{\mathrm{Pr}}


\usepackage{listings}
\newcommand{\passthrough}[1]{#1}
\lstset{% general command to set parameter(s)
  basicstyle=\small\ttfamily
  }


\usepackage[style=authoryear-comp,indexing=cite]{biblatex}
\ExecuteBibliographyOptions{sortlocale=general,maxnames=12,maxcitenames=3,parentracker=true,sortcites=false,uniquename=false,dashed=false,hyperref=true}
\uspunctuation
\addbibresource{references.bib}
\renewbibmacro{in:}{}


\begin{document}\sloppy

\title{\bf A Statistical Case for Qualified Scientific Optimism\vspace{.1in}}

\author{
Erik van Zwet\thanks{Department of Biomedical Data Sciences, Leiden University Medical Center.}
\and Andrew Gelman\thanks{Department of Statistics and Department of Political Science, Columbia University, New York.}
\and Witold Wi\k{e}cek\thanks{Development Innovation Lab, University of Chicago.}
\vspace{.1in}
}

\date{2 Apr 2026}

\maketitle

\begin{abstract}

We create an open-source database from empirical research in medicine, economics, psychology, psychotherapy, ecology, evolution, political science, education, and more, and make it available at \url{https://github.com/wwiecek/BEAR/}. Each result is represented by a $z$-value: an estimated effect divided by its standard error.

We then provide a modeling framework that allows us to move away from the usual attempts to classify effects as exactly zero or not, and instead focus on identifying the direction of effects and their replicability. 

After adjusting for publication bias, we find two patterns that hold across all datasets and their various subsets: (1) statistical significance alone does not imply a high probability of successful replication, but (2) it does imply a high probability that the direction of the observed effect was correct. 

We conclude that replication rate is a poor measure of the quality of a field of research and argue that an attitude of global scientific pessimism based on low replication rates is misguided. Instead, we favor a more optimistic view of scientific research based on the ability to identify the direction of effects. This optimism is not naive about the many challenges science continues to face such as questionable research practices (QRPs), paper mills, fake AI papers, and such. 

\end{abstract}

\section{Studying the quality of empirical science}\label{introduction}

The replication crisis has made us aware that there's a wide range of quality in the life and behavioral sciences, even considering subsets such as particular research areas and papers published in top journals.  This makes sense; if we roughly measure the quality of a study by its signal-to-noise ratio (the true effect size divided by its standard error), then we have some control over the denominator (by taking more or better measurements) but less control over the numerator. It's the nature of research to be uncertain and it's the duty of researchers to pursue all avenues, studying effects which might turn out to be null or highly context-dependent.

A key turning point in the replication crisis was the recognition and refutation of the false intuition that you can tell whether a study is good or not (tell whether a result is strong or not) using a statistical significance threshold.  Beyond the problems with any sharp threshold, even seemingly clear distinctions such as, for example, $p < 0.01$ compared to $p > 0.2$ do not discriminate well between high and low signal-to-noise ratios,\footnote{The two-sided $p$-values of 0.01 and 0.2 correspond to absolute $z$-values of 2.58 and 1.28, respectively. These are about equally likely when the signal-to-noise ratio is $(2.58+1.28)/2=1.93$. This signal-to-noise ratio corresponds to a little under 50\% power, as $\Pr(|z| \geq 1.96 \mid \text{SNR}=1.93) = 0.49$.} and a similar problem arises not just with $p$-values but also with Bayes factors and other measures of quality of evidence. 

But this should not be taken as a negative verdict on statistically-based science as a whole. To study this, we have created an extensive open-source database of empirical research results and focus here on 16 publicly available corpora in diverse fields.  Our optimistic finding is that, in these corpora, most results seem to be correct in sign.  Our remaining pessimism stems from clear signs of selection in analysis and reporting of results in some of the corpora, which points toward the need for real replication---not in the hope of confirming significance but to understand effects, and to minimize selection in all phases of the process.

It is well known to most that you shouldn't rely on a single study; replication and, ultimately, meta-analysis are the way to go.  But then many people are still dismayed when results ``fail to replicate".  One point of the present paper is that apparent replication rates in a research area can be low even when most studies are finding true effects. Hence apparent low replication rates should not be taken as implying that most of science is wrong.

\subsection{Going beyond the characterization of scientific claims as true or false}

In an influential paper, \textcite{ioannidis2005most} viewed empirical research as ``a $2 \times 2$ table in which research findings are compared against the gold standard of true relationships in a scientific field.''
%Assuming statistical tests are done at significance level $\alpha=0.05$ (two-sided), Ioannidis characterizes a research field in terms of the ratio of the number of ``true relationships'' to ``no relationships'' $R$ and the power, $1 - \beta$, the probability of obtaining a statistically significant result conditional on there being a true relationship. He proceeds to derive expressions for the $2 \times 2$ table of proportions of true and false positive and negative findings, which is also sometimes called the confusion matrix. Among these proportions, he emphasizes the  positive predictive value (PPV), the post-study probability of a research finding being true, conditional on its estimate achieving formal statistical significance. He finds $\mathit{PPV} = (1 - \beta)R/(R - \beta R + 0.05)$ and concludes that a research finding is more likely true than false if $(1 - \beta)R > 0.05$.

This $2 \times 2$ model is implicitly accepted in the meta-scientific literature, which tends to focus on statistical significance and replication rates \parencite{simmonsFalsePositivePsychologyUndisclosed2011, opensciencecollaborationEstimatingReproducibilityPsychological2015a, camerer2016evaluating}. In particular, low replication rates are taken by many as evidence of published empirical research being full of null claims and hence not credible as a whole, a position with political influence; see, for example, \textcite{garisto2025trump}.

We understand that the $2 \times 2$ model was always intended as an instructive simplification, introduced for the purpose of explaining that in an area of research where the power is low and there are few true relationships, many statistically significant results may not be true. However, far-reaching conclusions from this model have been made about the actual state of scientific research, in particular the claim that more than half of all published research findings are false.

We have two main criticisms of the $2 \times 2$ model. First, we disagree with the notion that the primary goal of scientific study is to separate ``no relationships'' or nulls from ``true relationships.''  Effects of interest are rarely exactly zero, and even if they are, they will no longer be zero if one takes leakage, expectation effects, dropout, etc. into account. But apart from that, with finite data, it is simply not possible to distinguish effects that are zero from those that are small.

Our second criticism is that the $2 \times 2$ model is not a realistic representation of empirical research. The key feature of the $2 \times 2$ model is that the PoS (or equivalently, the absolute value of the SNR) of the studies in some area of research take only two values; $\PoS=0.05$ ($\SNR=0$) for the null studies and some other value, for example $\PoS=0.8$ ($|\SNR|=2.8$), for the non-null studies. That is simply not possible. All researchers would have to know the true effect (if there is an effect) to be able to design studies to have either $\PoS=0.05$ or $\PoS=0.8$.

Here we make an effort to build more realistic statistical models of scientific research. We have compiled a database of corpora from medicine, economics, psychology, psychotherapy, ecology, evolution, political science, education, and more. Our models allow us to draw some tentative conclusions about concepts such as bias, coverage and replication.

In this paper we make two contributions. First, we make a large number of results from empirical research easily accessible in a new database, Benchmarks of Empirical Accuracy in Research (BEAR). We do not contribute any new data, but we process the available data to make them ready for meta-scientific research. We provide clear documentation and all data are available under open source license at \url{https://github.com/wwiecek/BEAR/}. We will continue to add new datasets as they become available.

Second, we extend the meta-scientific model of \textcite{stephens2017false,van2021statistical,zwet2022proposal}  to account for truncation and publication bias. This model allows us to study measures of research quality that are insensitive to whether effects are exactly zero or merely very small, such as the probability of correctly identifying the direction of an effect and the probability of replication \parencite{greenland2017invited, gelman2014beyond, gelman2000type}. 

\subsection{Summary of findings}

We estimate from our analysis that statistical significance usually implies a high probability that the direction of the observed effect is correct, but that it does not imply a high probability that a hypothetical exact replication study would reach statistical significance. In other words, it is expected that true findings (in the sense of correctly identifying the direction of an effect) will often not replicate. Similar conclusions were reached by \textcite{bak2022revisiting} and \textcite{neves2022most} on the basis of simulations and data from replication studies.

The scientific process is far from perfect, and questionable research practices, paper mills, and outright fraud are serious concerns \parencite{aquarius2025tackling,matusz2025threat}. But the mere fact that many studies have low probability of reaching statistical significance and therefore fail to replicate does not imply that most published research findings are false. 

The challenge is internalizing the juxtaposition of two ideas.  On one hand, there really is a replication crisis:  substantive and statistical theory, empirical evidence, and sociological reasoning all point toward the conclusion that traditional procedures for scientific quality control (peer review, causal identification, and statistical significance) are not enough to ensure the reliability of research findings.  On the other hand, the statistical evidence implies that in many fields, published results that reach $p<0.05$ generally fall in the correct direction.

%Evidence for the replication crisis:
%\begin{itemize}
%\item Substantive theory:  Many heavily-promoted claims (extra-sensory perception, power pose, mind-body healing, shark attacks, etc.) defy common sense.  This is not to say that they \emph{can't} be true, just that there is legitimate skepticism about them. Also, as discussed in the piranha paper, these claims can't all coexist.
%\item Statistics:  Researcher degrees of freedom, forking paths, type M and S errors, reasons why we can't take statistically significant findings at face value.
%\item Empirical evidence:  Various individual studies that did not replicate (elderly walking, etc.), also large-scale replication studies.
%\item Sociological reasoning:  Incentives for big claims, the feedback loop of overestimates and overconfidence.
%\end{itemize}
%A key takeaway from the replication crisis that we shouldn't automatically trust published results, even if they come from a respected institution and are published in a top journal.


\section{Data}

We have compiled a large number of empirical results from a variety of sources in a new database which we call Benchmarks of Empirical Accuracy in Research (BEAR). BEAR currently has 26 dataset entries, but for this paper we use 16 of those, comprising about 180,000 studies. Characteristics of the presented datasets are given in Table~\ref{tab:dataset_summary} and in Appendix \ref{app:datasets}. We describe our processing steps in Appendix \ref{standard-procedure-for-dealing-with-p-values-and-confidence-intervals}.

We distinguish three groups among our datasets. The first group contains eight curated sets. These include meta-analytic or systematic-review collections in intelligence research \parencite{nuijten2020intelligence}, economics \parencite{askarov2023significance}, psychotherapy \parencite{cuijpers2024effects}, political science \parencite{arel2022quantitative}, ecology and evolution \parencite{costello2022decline,yang2024large}, psychology \parencite{rodriguez2022psymetadata}, and medicine through the Cochrane Database of Systematic Reviews \parencite{schwab2024cochrane,cochrane2025cdsr}. We also include education intervention results from the What Works Clearinghouse \parencite{wwc2025studyfindings}. Where the source is meta-analytic, BEAR treats the rows as collections of individual studies, while retaining meta-analysis identifiers for other uses.

The second group consists of six corpora that can be thought of as representative samples of studies or results within research domains: SCORE dataset of claims in the social and behavioural sciences \parencite{tyner2026investigating}, main results from top economics papers \parencite{lang2025credible}, text-mined $t$-tests in cognitive neuroscience \parencite{szucs2017empirical}, clinical trial registry results from \texttt{clinicaltrials.gov} and EU CTR \parencite{clinicaltrials2025, euctr2025}, and text-mined $p$-values \parencite{chavalarias2016evolution} and confidence intervals \parencite{barnett2019examination} in biomedical journals.

For comparison, we also include a third group which has the studies that were studies that were selected for replication in two notable large scale replication efforts: Open Science Collaboration replications in psychology \parencite{opensciencecollaborationEstimatingReproducibilityPsychological2015a} and SCORE replications in the social and behavioural sciences \parencite{tyner2026investigating}. We will discuss these sets of studies together with the second group because they have very similar properties. 

\begin{table}
\centering
%\begin{tabular}{p{3.5cm}p{7.1cm}rrc}
\begin{tabular}{p{3.7cm}@{\hspace{1cm}}p{7cm}rrc}
Corpus & Purpose of dataset & studies & $\bar{k}$ & \% signif. \\
\toprule
\textcite{cochrane2025cdsr} & Database of systematic reviews in health and medicine & 30,306 & 1.3 & 31\%\\

ClinicalTrials.gov + EU CTR & Registries of clinical trials in the United States and European Union %\parencite{clinicaltrials2025,euctr2025}
& 24,468 & 2.0 & 52\%\\

\textcite{costello2022decline} & Do effects in ecology decline over time? & 12,927 & 6.8 & 42\%\\

\textcite{arel2022quantitative} & Assess statistical power in political science research & 2,252 & 7.4 & 47\%\\

\textcite{bartos2025exercise} & Effect of physical exercise on cognition, memory, and executive function & 2,239 & 1.0 & 26\%\\

\textcite{askarov2023significance} & Impact of mandatory data sharing on statistical significance in economics papers & 1,913 & 11.2 & 52\%\\

\textcite{nuijten2020intelligence} & Meta-meta-analysis of studies in intelligence research & 1,913 & 1.3 & 53\%\\

Metapsy (2025) & Database and meta-analyses of psychotherapy RCTs & 1,494 & 2.9 & 48\%\\

What Works Clearinghouse & Database of education intervention effect sizes \parencite{wwc2025studyfindings} & 1,408 & 8.6 & 34\%\\

psymetadata  & Curated psychology meta-analysis datasets \parencite{rodriguez2022psymetadata} & 721 & 11.8 & 36\%\\

\textcite{brodeur2024preregistration} & Pre-registration and pre-analysis plans in economics journals? & 176 & 47.9 & 37\%\\

\addlinespace
\midrule

Many Labs 2 \parencite{klein2018manylabs}  & Replication of classic and contemporary psychology experiments across labs & 1,414 & 1.0 & 44\%\\
\textcite{opensciencecollaborationEstimatingReproducibilityPsychological2015a} & Replications of a quasi-random sample of 100 experiments in psychology & 99 & 1.0 & 35\%\\

\addlinespace
\midrule

\textcite{chavalarias2016evolution} & Reporting of $p$-values over time & 1,887,178 & 4.2 & 84\%\\

\textcite{barnett2019examination} & Bias for statistical significance in health and medical journals & 416,027 & 3.1 & 83\%\\

\bottomrule
\end{tabular}
\caption{Datasets currently included in BEAR, divided into three groups: curated datasets, meta-analyses, and datasets based on large-scale scraping of published results; $\bar{k}$ is the mean number of observations per study, and ``\% signif.'' is the percentage of absolute $z$-values exceeding 1.96.}
\label{tab:dataset_summary}
\end{table}


While we do some data processing of our own, all of data collection and the vast majority of data processing was already done by the creators of the datasets we reference. Most of them already include $z$-values or effect sizes and standard errors, but in some cases we make calculations ourselves (for example, from $p$-values or counts for binary data); calculation details are in Appendix \ref{app:datasets}. 

Note that we can only identify the main claim in each study in two sets: \textcite{lang2025credible} and clinical trial registries. 
In most cases we have one or several estimates per study, see Table \ref{tab:dataset_summary}, but it can be as high as 50 $z$-values per study in the case of \textcite{arel2022quantitative}, which aimed to collect all tests reported by economics papers. Data are also heterogeneous in terms of fraction of $|z|$-values that exceed 1.96: from 0.31 in case of the Cochrane database, up to almost 90\% in SCORE and PubMed surveys.

% Our consolidated database has, for each published estimate, the data source, $z$-value, $p$-value, effect size estimate, standard error, sample sizes, year of publication, study type (e.g., randomized trial, observational study), measure type (e.g., difference of means, log risk ratio), and some categorization into fields of study and subgroups (e.g., medical field, clinical study phase). The current version of BEAR is a single table with xx rows. 

% In some cases we have access to summary data from which we can calculate effect estimates ourselves, but in most datasets we do not know the underlying method of calculation. Where only $p$-values are available, we assume they are two-sided and calculate the $z$-value as $z = \Phi^{-1}(1 - p/2)$. In four of the datasets, the $p$-values or $z$-values are truncated if researchers reported only $p < 0.05$, $p < 0.001$, etc. We account for this in our statistical model. 

\section{Statistical model}

We extend the model of \textcite{stephens2017false,van2021statistical,zwet2022proposal} to account for the truncation and publication bias which are conspicuously present in some of our datasets. The main goal is to estimate a latent distribution of the signal-to-noise ratios across a field of research and then use this distribution to derive various aspects of the sampling distribution of replication results.

\subsection{Signal-to-noise ratios}
 
Suppose that we have a collection of unbiased effect estimates which are normally distributed with known standard errors. The $z$-value is the ratio of the effect estimate to its standard error. We define the signal-to-noise ratio (SNR) as the ratio of the (unobserved) true effect to the standard error of its estimate. Our assumptions imply that the $z$-value is equal to the SNR plus an independent standard normal error. 

The null hypothesis that there is no effect is equivalent to the hypothesis that the SNR is zero. This hypothesis is commonly rejected when the absolute value of the $z$-value exceeds 1.96 or, equivalently, when the associated two-sided $p$-value is less than 0.05. The type I error probability of this test is $\alpha=0.05$. In our notation, 
$$\Pr(|z| \geq 1.96 \mid \text{SNR}=0)=0.05.$$ 

The probability of statistical significance (PoS) depends on the true effect and the standard error through the magnitude of the SNR,

\begin{equation}
\text{PoS}(\text{SNR})=\Pr(|z| \geq 1.96 \mid \text{SNR}) = \Phi(-1.96 - |\text{SNR}|) + 1 - \Phi(1.96 - |\text{SNR}|). \label{eq:power}
\end{equation}
For example, if the $\text{SNR}=2.8$ (or $-2.8$), then the PoS is 80\%. Averaging the PoS over the distribution of the SNR in a field of research, we obtain the ``assurance''  which we denote by $\overline{\text{PoS}}$. It is the proportion of studies for that field that reaches the 5\% level of statistical significance.  We can also consider  $\text{Pr}(|\text{SNR}| \geq 2.8) = \text{Pr}(\text{PoS(SNR)} \geq 0.8)$, which is the proportion of studies for which the power relative to the true effect size is 80\% or more.

The inherent variability of $p$-values can be most easily understood by considering SNRs and $z$-values. Suppose a study has 80\% power. In that case, the sampling distribution of the $z$-value is normal with mean 2.8 and unit variance. So, both $z=0.8 \ (p=0.42)$ and $z=4.8 \ (p=1.6 \cdot 10^{-6})$ can easily occur. 

\subsection{A model for the distribution of signal-to-noise ratios in a corpus}

We will use maximum likelihood to estimate the distribution of SNRs given a sample of absolute $z$-values. Our assumptions imply that $z$-values are equal to SNRs plus independent standard normal errors. We can therefore obtain the distribution of the absolute SNRs by deconvolution of the standard normal error component from the distribution of the absolute $z$-values \parencite{efron2016empirical,stephens2017false}. 

We construct the likelihood in four stages. First, we introduce a mixture distribution for the absolute $z$-values. Then, we add a parameter to account for publication bias. Next, we account for the fact that reported values are sometimes truncated. Finally, we weight the observations according to the number of $z$-values per study.

Our main assumption is that the distribution of the absolute $z$-values is well represented by a mixture of half-normal distributions. We will informally verify this assumption by inspection of the observed histograms. Noting that this assumption is equivalent to modeling the signed $z$-values as a mixture of zero-mean normal distributions, we have the mixture density,
\begin{equation}
f(z \mid p,\sigma) = \sum_{i=1}^{k} p_i\,\frac{1}{\sigma_i}
      \,\phi\!\left(\frac{z}{\sigma_i}\right),                  \label{eq:mixture}
\end{equation}
where $\phi(\cdot)$ is the standard normal density,   $p=(p_1,\dots,p_k)$ are the non-negative mixture weights that sum to unity and  $\sigma=(\sigma_1,\dots,\sigma_k)$ are the standard deviations. These standard deviations are at least 1 because we know that $z$-values have a standard normal noise component. We will use $k=4$, as we have verified that larger values of $k$ yield nearly identical results. 

The deconvolution to obtain the distribution of the SNRs is easy. We simply subtract 1 from the variances of the mixture components. So,
\begin{equation}
f_\text{SNR}(x \mid p,\sigma) = f(x \mid p,\sqrt{\sigma^2 -1}).
\end{equation}

To account for publication bias, we introduce a simple selection component in the likelihood. We follow \textcite{hedges1984estimation,hedges1992modeling} and define $\omega$ as the relative risk that a result with $|z|<1.96$ is observed compared to a result with $|z| \geq 1.96$
\begin{equation}
\omega \;=\;
\frac{\Pr(\text{result is published} \mid |z| <1.96)}
     {\Pr(\text{result is published} \mid |z| \ge1.96)}.
\end{equation}
\noindent
The probability of publication is
\begin{equation}
C(p,\sigma,\omega) = \omega\,\Pr(|z| < 1.96 \mid p,\sigma) + \Pr(|z| \geq 1.96 \mid p,\sigma)
\end{equation}
and the density of the $z$-values conditional on publication becomes

\begin{equation}
f(z \mid p,\sigma, \omega, \text{pub}) = 
\begin{cases}
\omega\,f(z \mid p,\sigma) / C(p,\sigma,\omega) \quad   & |z| \leq 1.96 \\
      f(z \mid p,\sigma)/ C(p,\sigma,\omega)                              & \text{otherwise.}
\end{cases}
\end{equation}

The $z$-value is sometimes left or right truncated, most often at 1.96, but also at, for example, 2.57 (when papers report $p < 0.01$), 3.29 ($p < 0.001$), or 1.64 ($p < 0.1$). To take this into account, we introduce a censoring indicator $\delta$ taking the values $<$, $>$, or $=$ for left censoring, right censoring, or no censoring, respectively. We denote the (possibly) censored value of the $z$-value by $\tilde{z}$. The distribution of the pair $(\delta,\tilde{z})$ conditional on publication is
\begin{align*}
&f(\tilde{z},\delta \mid p,\sigma, \omega, \text{pub})= \\
&\begin{cases}
\omega\,f(\tilde{z}\mid p,\sigma)/ C(p,\sigma,\omega)         & \delta=\texttt{"="},\; \tilde{z}<1.96,\\[4pt]
      f(\tilde{z}\mid p,\sigma)/C(p,\sigma,\omega)              & \delta=\texttt{"="},\; \tilde{z}\ge 1.96,\\[6pt]
\omega\,F(\tilde{z}\mid p,\sigma)/ C(p,\sigma,\omega)            & \delta=\texttt{"<"},\; \tilde{z}< 1.96,\\[4pt]
(\omega\,F(1.96\mid p,\sigma)+F(\tilde{z}\mid p,\sigma)-F(1.96\mid p,\sigma))/ C(p,\sigma,\omega)          & \delta=\texttt{"<"},\; \tilde{z}\ge  1.96,\\[6pt]
(1-F(\tilde{z}\mid p,\sigma))/ C(p,\sigma,\omega)          & \delta=\texttt{">"},\; \tilde{z}\ge 1.96,\\[4pt]
((1-F(1.96\mid p,\sigma))+\omega\,(F(1.96\mid p,\sigma)-F(\tilde{z}\mid p,\sigma)))/ C(p,\sigma,\omega)
                                               & \delta=\texttt{">"},\; \tilde{z}<  1.96.
\end{cases}
\end{align*}
\noindent
The first two rows cover the case without truncation. Rows three and four are left censoring (statements such as $p > 0.05$), and  rows five and six deal with right censoring (e.g., $p < 0.001$). We also introduce artificial censoring when $z$-values are reported to be exactly zero, by treating them as $|z| < 0.5$.

Most of the datasets have multiple $z$-values per study. To prevent studies with many $z$-values from dominating the likelihood, we weight observations by $w_j$, the inverse of the number of $z$-values originating from study $j$. This does not account for the dependence among $z$-values from the same study but that does not concern us because we do not attempt to quantify the uncertainty of our estimates.  The reason is that our datasets are so large that the sampling uncertainty which determines standard errors and confidence intervals, is negligible compared to model uncertainty. The full weighted log likelihood is  
\begin{equation}
\ell(p,\sigma, \omega)
= \sum_{j=1}^{n} w_j \log f(\tilde{z_j},\delta_j \mid p,\sigma, \omega, \text{pub}).               \label{eq:loglik}
\end{equation}
\noindent
Even for large datasets, numerical optimization of this log likelihood is feasible, taking just minutes.  For the paper fits, we cap the fitting input at 50,000 rows per corpus, using study-level thinning as described in the Data section.

The estimated distribution of the $z$-values without publication bias is $f(z \mid \hat p,\hat \sigma)$ as defined in (\ref{eq:mixture}) and consequently the estimated distribution of the SNRs is
\begin{equation}
f_\text{SNR}(x \mid \hat p,\hat \sigma) = f(x \mid \hat p,\sqrt{\hat \sigma^2 -1}).
\end{equation}

We can derive many important quantities from the distribution of the SNRs. For example, the PoS is a function of the SNR, so we can obtain its distribution by transforming the distribution of the SNRs. In particular, the proportion of studies that have at least 80\% probability of reaching statistical significance is
\begin{equation}
\Pr(\text{PoS} \geq 0.8) = \Pr(|\text{SNR}| \geq 2.8).
\end{equation}
The marginal density of the SNRs, together with the conditional density of the $z$-values given the SNRs, determines the joint distribution. And from the joint  distribution we can derive the conditional distribution of the SNR given the $z$-value.
We provide the formulas in Appendix \ref{appx_details}. In turn, this allows us to derive the conditional probability of the correct sign,
\begin{equation}\label{eq:sign}
\text{Probability of the correct sign} = \Pr(z \cdot \text{SNR} > 0 \,\big|\, |z|),
\end{equation}
Finally, we imagine that a second replication study with exactly the same SNR yields another $z$-value $z_\text{repl}$. This  $z_\text{repl}$ is conditionally independent of the original $z$ given SNR, and has the same distribution. Thus \textcite{vanzwet2022large} were able to derive the conditional distribution of $|z_\text{repl}|$ given $|z|$. In particular, we obtain 
\begin{equation}\label{eq:repl}
\text{Probability of ``successful replication''} = \Pr(z_\text{repl} \cdot z > 0 \text{ and } |z_\text{repl}| > 1.96 \,\big|\, |z|) 
\end{equation}

All the computations that yielded the results reported in the next section are documented in the online supplement to this paper. 

\section{Results}

Figure~\ref{fig:zmixtures} displays the histograms of the observed absolute $z$-statistics of our 14 non-replication corpora. The solid curves are the fitted mixtures of half-normal distributions conditional on publication, and the dashed curves are the ``corrected'' unconditional distributions. We find that the estimated distributions generally track the histogram well and conclude that our model is appropriate for our purpose of quantifying large-scale properties of research areas. The first eight panels are curated sets and the next six are representative domain samples, arranged within groups by their estimated assurances $\overline{\text{PoS}}$. Figure~\ref{fig:replication_mixtures} compares the original-study and replication-study distributions for SCORE and OSC. 

Mixtures of half-normal densities are necessarily decreasing. Apart from the effects of publication bias, we also see this feature in the histograms. We claim that this is typical of large collections of $z$-values. To further support this claim, we have broken down the Cochrane data by 19 medical specialties and the data from {\tt clinicaltrials.gov} by trial phase, see Figure~\ref{fig:subsets}. Again, we see decreasing distributions of absolute $z$-values.

The discontinuities (``jumps'') of the densities at $|z|=1.96$ are a feature of Hedges' selection model. It is tempting to focus on the extent of publication bias, which may in principle be quantified by the estimated relative risk of publication $\hat \omega$ as reported in Table~\ref{tab:results_summary_table2}. However, there are several reasons why we must resist that temptation. Many studies in our datasets have multiple outcomes, and we do not know which of them drive publication decisions. Moreover, while Hedges' single-parameter selection model is useful for adjusting our estimates of the distributions of the $z$-statistics, it seems too limited for reliably quantifying and comparing the extent of publication bias across fields of research. This is also not the goal of our study.

The main difference between the distributions of the $z$-values is their widths. A wider distribution of $z$-values implies a wider distribution of SNRs. In other words, the studies in that research area tend to have larger absolute SNRs. Since the PoS is an increasing function of the absolute SNR, this also means that a higher proportion of results will reach statistical significance. 

The two columns in Table~\ref{tab:results_summary_table2} that are labeled ``Significance'' report the observed proportions of significant results and the estimated (``corrected'') assurances in the 16 corpora. Recall that the assurance $\overline{\text{PoS}}$ is the proportion of significant results, adjusted for publication bias; see (\ref{eq:power}).

The assurance is strictly less than the raw proportion of significant results, unless $\hat{\omega}=1$ so there is no discernible publication bias. Across the 16 corpora, the estimated assurances range from 0.23 to 0.52.  The estimated assurance of phase 3 clinical trials in Figure~\ref{fig:subsets} matches the estimated success rates (probability of moving to regulatory approval following trial) in phase 3 trials in drug development reported by \textcite{thomas2021clinicaldevelopment}.

By transforming the distribution of the signal-to-noise ratios, we obtain the full distribution of the PoS within each dataset. The left panel of Figure~\ref{fig:pos_gap} displays the complementary cumulative distribution functions (survival functions) of the PoS in the 16 datasets. In other words, it shows what proportions of studies in each dataset ``survive'' to achieve a given level of power. We find that between 10\% and 33\% of studies achieve 80\% PoS; see Table~\ref{tab:results_summary_table3}. 

Since the great majority of studies are targeted at reaching 80\% power, the proportion of studies that reach 80\% PoS may be tentatively interpreted as the proportions of studies where the assumptions of the sample size calculation were true. In general, this proportion is low. However, sample size calculations are not supposed to be aimed at the true effect (which is never known) but at the effect one would not want to miss \parencite{senn2002cross}. In that sense, a low proportion of studies with 80\% PoS just means that it is hard to come up with interventions that have large effects. Pressures of time, money, and logistics also can drive down sample sizes.  

The three columns in Table~\ref{tab:results_summary_table2} that are labeled ``Successful replication'' report what would happen if a randomly selected study from a particular corpus would be replicated exactly. We show the probability of obtaining a statistically significant result in the same direction both unconditionally and conditionally on the original study having $|z|=1.96$ or $|z|\geq 1.96$; see equation (\ref{eq:repl}). We find that if a result is marginally statistically significant at the 5\% level (i.e., $|z|=1.96$) then the probability of ``successful replication'' is never much more than the base rate, and often even less. This will be a surprise to those who think of statistical significance as a stand-in for replication. This phenomenon is {\em not} a consequence of many studies having low SNRs (low PoS). If all the studies had higher SNRs, then the assurance would be higher but observing $|z|=1.96$ would be even less impressive.

The three columns that are labeled ``Correct sign'' report
the estimated proportion of studies where the observed effect has the correct direction (sign). Again, we also condition on $|z|=1.96$ or $|z|\geq 1.96$, cf as defined in equation (\ref{eq:sign}).

If we focus on the case where $|z|=1.96$ then we find a large difference between the probabilities of successful replication and the correct sign. Across the 16 corpora, the estimated probability of replication ranges from 27\% to 44\% while the estimated probability of correct sign ranges from 82\% to 97\%. 

The right panel of Figure~\ref{fig:pos_gap} displays the probabilities of successful replication and correct sign conditional on the observed absolute $z$-value, see also Figure~\ref{fig:gap_plot_all}.  The probability of correct sign is always much higher than the probability of successful replication.

\begin{figure}
\centering
\includegraphics[width=1\textwidth]{fig/mixtures_plot_collection_groups.pdf}
\caption{Histograms of absolute $z$-values across 14 non-replication corpora with fitted curves overlaid on top. The solid curves are densities, fitted using mixtures of half-normal distributions and conditioning on publication. Dashed curves are the unconditional ``corrected'' distributions, scaled to match the solid curve in region $|z| > 1.96$. Blue: curated sets; red: representative domain samples.}\label{fig:zmixtures}
\end{figure}

\begin{figure}
\centering
\includegraphics[width=1\textwidth]{fig/mixtures_plot_replication_projects.pdf}
\caption{Histograms of absolute $z$-values for the SCORE and OSC original-study and replication-study distributions. Only the unconditional ``corrected'' mixture density is shown. Grey panels show original studies; yellow panels show replication studies.}\label{fig:replication_mixtures}
\end{figure}


\begin{table}
\centering
\scriptsize
\begin{tabular}{lccccccccc}
\multicolumn{2}{c}{}
& \multicolumn{2}{c}{Significance}
& \multicolumn{3}{c}{``Successful'' replication}
& \multicolumn{3}{c}{Correct sign} \\
\cmidrule(lr){3-4}\cmidrule(lr){5-7}\cmidrule(lr){8-10}
& & & & & \multicolumn{2}{c}{\scriptsize Conditional on}
& & \multicolumn{2}{c}{\scriptsize Conditional on} \\
Corpus & \multicolumn{1}{c}{\scriptsize $\hat\omega$}
        & \multicolumn{1}{c}{\scriptsize signif.}
        & \multicolumn{1}{c}{\scriptsize $\overline{PoS}$}
        & \multicolumn{1}{c}{\scriptsize uncond.}
        & \multicolumn{1}{c}{\scriptsize $|z| = 1.96$}
        & \multicolumn{1}{c}{\scriptsize $|z| \geq 1.96$}
        & \multicolumn{1}{c}{\scriptsize uncond.}
        & \multicolumn{1}{c}{\scriptsize $|z| = 1.96$}
        & \multicolumn{1}{c}{\scriptsize $|z| \geq 1.96$} \\
\cmidrule(lr){1-10}
\multicolumn{10}{l}{\emph{Curated sets}}\\
\cmidrule(lr){1-10}
Nuijten et al & 0.85 & 0.53 & 0.48 & 0.47 & 0.41 & 0.75 & 0.88 & 0.96 & 0.99\\
Askarov et al & 0.70 & 0.52 & 0.48 & 0.46 & 0.40 & 0.75 & 0.88 & 0.96 & 0.99\\
Metapsy & 0.73 & 0.48 & 0.43 & 0.41 & 0.40 & 0.70 & 0.87 & 0.96 & 0.99\\
Arel-Bundock et al & 0.64 & 0.47 & 0.40 & 0.38 & 0.36 & 0.68 & 0.85 & 0.96 & 0.99\\
Costello and Fox & 0.74 & 0.42 & 0.39 & 0.37 & 0.35 & 0.74 & 0.79 & 0.89 & 0.98\\
psymetadata & 0.66 & 0.36 & 0.32 & 0.30 & 0.30 & 0.74 & 0.72 & 0.81 & 0.96\\
What Works Clearing. & 0.88 & 0.34 & 0.32 & 0.30 & 0.31 & 0.69 & 0.75 & 0.87 & 0.97\\
Cochrane & 0.82 & 0.33 & 0.28 & 0.26 & 0.30 & 0.60 & 0.77 & 0.90 & 0.97\\
\addlinespace
\cmidrule(lr){1-10}
\multicolumn{10}{l}{\emph{Representative domain samples}}\\
\cmidrule(lr){1-10}
Chavalarias et al & 0.15 & 0.87 & 0.52 & 0.50 & 0.44 & 0.76 & 0.89 & 0.97 & 1.00\\
SCORE, all claims & 0.12 & 0.87 & 0.47 & 0.45 & 0.41 & 0.73 & 0.88 & 0.96 & 0.99\\
Lang & 0.24 & 0.58 & 0.42 & 0.41 & 0.38 & 0.70 & 0.86 & 0.96 & 0.99\\
Szucs and Ioannidis & 0.35 & 0.62 & 0.35 & 0.33 & 0.36 & 0.62 & 0.84 & 0.96 & 0.99\\
clinicaltrials.gov + EU CTR & 0.56 & 0.44 & 0.33 & 0.32 & 0.30 & 0.73 & 0.74 & 0.85 & 0.97\\
Barnett and Wren & 0.08 & 0.84 & 0.30 & 0.28 & 0.30 & 0.58 & 0.81 & 0.94 & 0.98\\
\addlinespace
\cmidrule(lr){1-10}
\multicolumn{10}{l}{\emph{Replication efforts}}\\
\cmidrule(lr){1-10}
SCORE replications & 0.70 & 0.64 & 0.51 & 0.49 & 0.39 & 0.78 & 0.88 & 0.96 & 0.99\\
OpenSciCollab replications & 1.05 & 0.35 & 0.36 & 0.35 & 0.31 & 0.74 & 0.80 & 0.91 & 0.98\\
\bottomrule
\end{tabular}
\caption{Summary of signal-to-noise ratio modeling results for the 16 corpora used in the paper. Rows are grouped into curated sets, representative domain samples, and replication projects. Calculations of power, replication, and direction of effects use the distributions of $z$-values and signal-to-noise ratios by fitted mixture models without selection.}
\label{tab:results_summary_table2}
\end{table}


\begin{figure}
\centering
\includegraphics[width=.9\textwidth]{fig/pos_gap.pdf}
\vspace{-.2in}
\caption{Left panel: ``survival curves'' of the probability of significance: $1 - (\mbox{CDF of PoS})$ for each of the 16 corpora. Right panel: conditional probabilities that the observed direction of the effect matches the true direction (top set of curves) and the conditional probability of ``successful'' replication (bottom set). Colors distinguish curated sets, representative domain samples, and replication efforts.}\label{fig:pos_gap}
\end{figure}

The estimated relative probability of reporting, $\hat \omega$, can be as low as 0.08. However, Figure~\ref{fig:zmixtures} highlights a clear divide: we find much less selection on significance in the curated datasets, with $\hat \omega > 0.66$ across all eight sets, than in the domain-representative samples, where $\hat \omega$ ranges from 0.08 for \textcite{barnett2019examination} to 0.50 for clinical trial registries. However, we must be careful not to over-interpret this result. The corpus of \textcite{barnett2019examination} is based on confidence intervals, and it is possible that non-significant results were actually reported, just not with the associated confidence intervals. 


\section{Discussion}

\subsection{An alternative to the binary model of scientific discovery}

The  $2 \times 2$ model \parencite{ioannidis2005most} is often taken to be a reasonable, if simplified, description of scientific research. According to this model, the probabilities of significance (PoS) of the ``true relationships'' and ``no relationships'' are well separated. The PoS of ``no relationships'' is 5\% by the design of the statistical test, and the PoS of ``relationships'' is much larger by the design of the experiment. If the ratio of ``true relationships'' to ``no relationships'' is not too low, then we can reliably separate them. That is, the positive predictive value (the conditional probability of a true relationship given statistical significance) and the negative predictive value (the conditional probability of no relationship given non-significance) are both high. Moreover, the conditional probability of successful replication (reaching statistical significance again in an exact replication study) is also high.

We have developed an alternative statistical model that goes beyond the $2\times 2$ framework by allowing the PoS to range continuously from 0.05 to 1 and uses Hedges' model to adjust for selective reporting. We also allow for truncation. We fit our model to sixteen corpora across diverse fields collected from publicly available sources.

The model fits well; see Figures~\ref{fig:zmixtures} and \ref{fig:subsets}. In contrast, the $2 \times 2$ model does not fit our data, as we observe no clear separation between the probability of statistical significance for ``true relationships'' and ``no relationships.''



\subsection{What our meta-analysis says about effect sizes and replication rates}

There are four notable features that hold across all the datasets of empirical research we collected. 

First, the absolute $z$-values have a decreasing density, which suggests that the absolute SNRs also have a decreasing density. In other words, studies with low SNRs are more common than those with high SNRs. We do find that there are considerable differences between research domains in the widths of these distributions which means that the studies in some areas tend to have higher SNRs than in other areas. 

Second, the probability of significance (PoS) tends to be low and only a small minority of studies achieve 80\% power. 

Third, reaching statistical significance does not imply a high probability of replication, see Figure~\ref{fig:pos_gap} and Figure~\ref{fig:gap_plot_all}. Even at $p = 0.001$ ($|z|=3.29$) replication is far from assured. 

Fourth, statistical significance {\em does} imply that the sign of the effect is likely correct. 

\subsection{Understanding the literature on meta-replications}

A low replication rate does not mean that studies in a field can not be trusted.  As shown in Table \ref{fig:pos_gap}, even in high-quality corpora expected probabilities of successful replication are below 50\%, as low as 20\%. But we still estimate that over 70-90\% of estimates have the correct sign. We believe that replication studies are valuable as long as these two statistical properties are understood.  

We are not saying that any study that appears in print has a high replication probability; however, if an experiment comparable to those in our corpus is done on a new topic, and it reaches a preregistered $z$-score in the range of 2 to 3, we would expect roughly a 90\% chance of getting the correct sign and a 50\% probability of statistically significant replication, as shown in Figure \ref{fig:pos_gap}b.

However, a lot of research is still not preregistered and we should be concerned about researcher degrees of freedom and forking paths. Our model does not quantify this aspect of researchers' behavior. However, we note a useful, albeit pessimistic, heuristic for judging performance of an average study, which is to ignore claims of statistical significance. For example, \textcite{opensciencecollaborationEstimatingReproducibilityPsychological2015a} reported success with 36 out of 100 replication studies in psychology. 97 out of 100 original results were statistically significant. Per Table \ref{tab:results_summary_table2} we find that conditioning on $|z| > 1.96$ we would expect higher replication rates, 58-78\%, in all corpora, including in the replication project. On the other hand, the observed replication rate is close to the unconditional replication probability in Table \ref{tab:results_summary_table2}.

Our interpretation of \textcite{miske2026,tyner2026} . . . 

\subsection{Understanding our findings using the anthropic principle}

Our approach is meta-scientific in that it is based on a global analysis of large corpora of published results. 
Why do these different corpora, coming from many different experiments in different fields, show similar distributions of the SNR? We tentatively ascribe this to a sort of anthropic principle \parencite{anthropic} whereby researchers respond to the pressures of time, money, and logistics on one side and the goal to reach statistical significance on the other, by making their studies just about large enough to assess the direction of the effect of interest, but not larger. 

That said, these incentives do not apply in all areas, and there are subfields of science that are plagued by questionable research practices and increasingly influenced by paper mills and fraudulent research. 
If we have good prior reasons to doubt the study design or integrity of the analysis, we should use that information. In other words, not every result should be treated as an typical draw from the broad reference class. Our qualified optimism extends only to research for which we do not have decisive reasons to distrust.


\subsection{Reasons for qualified optimism about empirical research}

% WW: I think this is true, but it's a point also made elsewhere in the paper:
% Instead of focusing on the paradigm of ``no relationships'' versus ``true relationships'' and associated type I and type II errors, our model allows us to examine replication and identifying the direction of effects which are both insensitive to whether effects are zero or merely small. 

% WW: This is a repeat of previous subsection
%We find that statistical significance usually implies a high probability that the direction of the observed effect is correct, but that it does not imply a high probability that an exact replication study will be significant again.

The standard practice in empirical research is that if a result reaches a statistical significance threshold, it is treated as real and the corresponding estimate is taken at face value. Results that are not significant are presented as a null findings or, at best, scenarios where there is not enough information to draw any conclusions. Indeed, it can be viewed as scandalous when a statistically significant published result does not replicate \parencite{tversky1971}. 

One way to summarize our result is to say that empirical research often leads to true findings---in the sense of correctly identifying the direction of an effect---that do not replicate in the sense of achieving statistically significant result in the same direction as the original study.  Therefore, failure to replicate should not be taken to mean that the original result was a fluke or a fake. The pessimism in much of the meta-scientific literature reflects a one-sided interpretation of Figure~\ref{fig:pos_gap} and Figure~\ref{fig:gap_plot_all} focusing only on the low replication rates. This skewed view has serious consequences as it may lead to dangerous nihilism among policy makers \parencite{garisto2025trump}. With our study, we hope to rectify it.

%WW: I think these next two paragraphs are very good points, but not sure where best to fit them and I think the current flow is good
%This should not come as a surprise, as the difference between ``significant'' and ``not significant'' is not itself statistically significant \parencite{gelman2006difference}. 
%This suggests a widespread discrepancy between expectations and reality which we believe to be due to the unrealistic $2 \times 2$ model that people have in mind.

%EvZ: commented out becuase it's reptitive
%Along with our qualified optimism about the direction of effects, our analysis suggests that a low observed rate of statistically significant replications should not be taken as evidence that most published claims are false.  
%The bad news is that most studies have less than 80\% probability of significance (power relative to the true effect size); the good news is that, even though average probabilities of significance range from 20\% to 50\%, replications are likely to go in the right direction, with estimated average type S error rates in our corpora in the range of 1\%--3\%. 

\section{Acknowledgments}
We thank all the compilers of the corpora and the authors of the original studies for their great effort. In particular, we are grateful to Frantisek Bartos, Ryan Briggs, Chris Doucouliagos, Chris Mead, Martina Sladekova, and Tom Stanley for helping us understand the datasets (we do not mean to suggest that they necessarily agree with our conclusions).  We thank the U.S. Office of Naval Research for partial support of this work.

% This is now leading into the follow-up paper---thinking at the level of individual studies?

%In contrast, in our corpus-informed Bayesian approach, estimates are not taken at face value---they are partially pooled toward zero---even in the case of statistical significance, thus adjusting for possible type M (magnitude) errors \parencite{van2021statistical,zwet2022proposal}.  And, as discussed in the present article, even non-statistically-significant results are not taken to be zero, as estimates are still likely to have the right sign.  We estimate the average type S error rates in our corpora in the range 10\%--30\%. 

%We need to strike a balance . . .
%
%
%It depends on the subfield! . . .
%
%
%\section*{To keep in mind}
%
%{\bf Disclaimer}: Our results apply to repeated sampling of studies (or a random sample of a study) from certain research fields. To apply them to a particular study, would require us to ignore everything about that study except the observed $z$-value.
%
%\medskip\noindent
%{\bf Main message 1}: Empirical finding: research areas are similar in the sense that they have decreasing distributions of absolute $z$-values (and hence decreasing distributions of absolute SNRs). This is an anthropogenic/anthropic phenomenon.
%
%\medskip\noindent
%{\bf Main message 2}: Individual studies can be bad or untrustworthy for many reasons. However, a replication rate of 30-40\% of a research area is no reason for mistrust. It's merely a consequence of the inherent variability of $z$-values, i.e. they're the convolution of some distribution that is mostly supported on $[-2,2]$ (the SNRs) and the standard normal.
%
%\bigskip
%The posts on statmodeling about $z$ and SNR show that there’s a lot of confusion. The main issues seem to be:
%\begin{itemize}
%\item Confusing the SNR with the effect. Even if the effect of some intervention is large, the distribution of $|\SNR|$ can still be decreasing.
% 
%\item Confusing selection on large $z$ (bad) with selection on large SNR (good).
% 
%\item Confusing the distribution of z across some field of study with the conditional distribution of z given SNR. The former is not supposed to look normal. In the absence of selection on $z$, it’s only supposed to be as smooth as N(0,1). That the distribution of $|z|$ is often decreasing is an empirical finding.
%\end{itemize}
%We’ll need to work very hard to explain these things! Maybe think of some good visuals?
%

\printbibliography[title=References]

\pagebreak
\clearpage

\appendix
\renewcommand{\thefigure}{S\arabic{figure}}
\renewcommand{\thetable}{S\arabic{figure}}
\setcounter{figure}{0}
\section{Details of calculations}\label{appx_details}

Suppose the marginal distribution of the $z$-value is a mixture of zero-mean normal distributions,
\begin{equation*}
f(z \mid p,\sigma) = \sum_{i=1}^{k} p_i\,\frac{1}{\sigma_i}
      \,\phi\!\left(\frac{z}{\sigma_i}\right).
\end{equation*}
Suppose also that the conditional distribution of the $z$-value, given the SNR, is normal with mean SNR and variance 1. Together, the marginal and conditional distributions determine the joint distribution of the $z$-value and the SNR. Therefore, we also have the conditional distribution of the SNR given the $z$-value. It is a mixture of normal distributions with weights $q_i(z)$, means $\mu_i(z)$, standard deviations $\tau_i$ given by

\begin{equation}q_i(z) = \frac{p_i \phi(z/\sigma_i)/\sigma_i }{f(z \mid p,\sigma)}, \quad \mu_i(z)=\frac{\sigma_i^2 - 1}{\sigma_i^2}\,z \quad\text{and}\quad \tau_i=\frac{\sqrt{\sigma_i^2 - 1}}{\sigma_i}.\end{equation}

\noindent
The conditional probability of the correct sign is
\begin{equation}
\Pr(z \cdot \text{SNR} > 0 \mid |z|) = \sum_{i=1}^{k} q_i(z) \Phi\left( \frac{\mu_i(|z|)}{\tau_i}\right).
\end{equation}

If $z_\text{repl}$ is the $z$-value of an exact replication study, then the conditional distribution of $z_\text{repl}$ given the $z$-value of the original study is a mixture of normal distributions with weights $q_i(z)$, means $\mu_i(z)$ and variances $\tau_i^2 +1$. The conditional probability of ``successful replication'' is
\begin{equation}
\Pr(z_\text {repl} \cdot z > 0 \text{ and } |z_\text{repl}| > 1.96 \mid |z|)  = \sum_{i=1}^{k} q_i(z) \Phi\left( \frac{\mu_i(|z|) - 1.96}{\sqrt{\tau_i^2+1}}\right). 
\end{equation}

\begin{figure}[h]
\centering
\includegraphics[width=16.5cm]{fig/gap_plot_all.pdf}
\caption{Relationship between absolute $z$-values and the probability that the direction of the observed effect matches the true direction (blue curves) and the probability of ``successful replication'' (red curves). The vertical line is $\overline{\mathrm{PoS}}$. }
\label{fig:gap_plot_all}
\end{figure}

\begin{figure}[h]
    \centering
    \includegraphics[width=1\linewidth]{fig/mixtures_plot_subsets.pdf}
    \vspace{-.2in}
    \caption{Subsets of interventional treatment studies from {\tt clinicaltrials.gov}, broken down by phase (top row) and from Cochrane database (bottom row), grouped by specialty (out of 19 included in the dataset): two with lowest $\overline{\mathrm{PoS}}$ on the left and two with the highest on the right. The peaks in the histograms at $|z| = 1.96, 2.58, 3.29$ are due to trials reporting (two-sided) $p<0.05, p<0.01, p<0.001$.}
    \label{fig:subsets}
\end{figure}

\begin{table}
\centering
\begin{tabular}{lcccccc}
\multicolumn{2}{c}{} 
& \multicolumn{2}{c}{Significance}
& \multicolumn{3}{c}{PoS at least 80\%} \\
\cmidrule(lr){3-4}\cmidrule(lr){5-7}
& & & & & \multicolumn{2}{c}{\scriptsize Conditional on} \\
Corpus & \multicolumn{1}{c}{\scriptsize $\hat\omega$}
        & \multicolumn{1}{c}{\scriptsize signif.}
        & \multicolumn{1}{c}{\scriptsize $\overline{PoS}$}
        & \multicolumn{1}{c}{\scriptsize uncond.} 
        & \multicolumn{1}{c}{\scriptsize $|z| = 1.96$}
        & \multicolumn{1}{c}{\scriptsize $|z| \geq 1.96$} \\
\cmidrule(lr){1-7}
Cochrane & 0.70 & 0.31 & 0.23 & 0.10 & 0.05 & 0.40\\
ctgov / EU CTR & 0.95 & 0.52 & 0.47 & 0.28 & 0.09 & 0.58\\
Costello and Fox & 0.79 & 0.42 & 0.39 & 0.23 & 0.08 & 0.58\\
Arel-Bundock et al. & 0.65 & 0.47 & 0.40 & 0.20 & 0.07 & 0.48\\
Bartos et al. & 0.82 & 0.26 & 0.22 & 0.08 & 0.03 & 0.35\\
Askarov et al. & 0.71 & 0.52 & 0.48 & 0.30 & 0.10 & 0.59\\
Nuijten et al. & 0.85 & 0.53 & 0.48 & 0.29 & 0.11 & 0.57\\
Metapsy & 0.83 & 0.48 & 0.46 & 0.26 & 0.10 & 0.54\\
What Works & 0.88 & 0.34 & 0.32 & 0.17 & 0.06 & 0.51\\
psymetadata & 0.67 & 0.36 & 0.33 & 0.21 & 0.08 & 0.62\\
Brodeur et al. & 0.79 & 0.37 & 0.36 & 0.17 & 0.05 & 0.46\\
\cmidrule(lr){1-7}
Many Labs 2 & 0.99 & 0.44 & 0.44 & 0.35 & 0.11 & 0.78\\
OpenSciCollab & 1.00 & 0.35 & 0.36 & 0.22 & 0.09 & 0.60\\
\cmidrule(lr){1-7}
Chavalarias et al. & 0.20 & 0.84 & 0.51 & 0.32 & 0.13 & 0.60\\
Barnett and Wren & 0.08 & 0.83 & 0.32 & 0.12 & 0.04 & 0.37\\
\cmidrule(lr){1-7}
\end{tabular}
\caption{Summary of signal-to-noise ratio modeling results for the 15 corpora of studies in our data. Calculations of power, replication, and direction of effects use the distributions of $z$-values and signal-to-noise ratios by fitted mixture models without selection.}
\label{tab:results_summary_table3}
\end{table}


\clearpage
\section{Datasets included in BEAR}\label{app:datasets}
% Do not edit by hand; generated from paper/datasets.Rmd.
This appendix describes the 16 corpora used in the paper. Note that the
full BEAR database contains additional datasets. More complete
documentation is available online.

\newpage

\subsection{psymetadata and Nuijten et al (psychology
datasets)}\label{psymetadata-and-nuijten-et-al-psychology-datasets}

\textbf{Reference:} A set of psychology datasets compiled by
\textcite{rodriguez2022psymetadata}. We also disaggregate a large
meta-meta-analysis of intelligence research
\autocite{nuijten2020intelligence}.

\textbf{Snapshot date:} January 2026

\textbf{Research question:} \emph{psymetadata} is an R package that
curates multiple open datasets from published meta-analyses in
psychology for teaching and demonstrations.
\textcite{nuijten2020intelligence} is a large meta-meta-analysis of
intelligence research, which focused on power and small study effects.

\textbf{Data availability:} 22 datasets are distributed with the
\texttt{psymetadata} R package (CRAN). We do not use its Many Labs 2
dataset, which we obtain from the original source and include
separately. After excluding Many Labs 2 and separately representing the
intelligence meta-meta-analysis, the remaining psychology datasets range
from 65 rows to over 1,000. The package is GPL-3; individual datasets
originate from the cited source papers.

\textbf{Data processing:} Processing was minimal. We used regular
expressions to extract year labels from titles. We dropped some invalid
rows (missing, non-finite, or non-positive standard errors). Signed
z-values were computed as effect sizes divided by their standard errors.

\textbf{Study characteristics:}

\begin{itemize}
\item
  \textbf{psymetadata:} the designs of the primary studies are not
  harmonised across the imported datasets.
\item
  \textbf{psymetadata:} the effect-size measure is not harmonised across
  the imported datasets.
\item
  \textbf{psymetadata:} records are grouped by original psymetadata
  dataset.
\item
  \textbf{Nuijten:} the designs of the primary studies are not
  classified.
\item
  \textbf{Nuijten:} BEAR retains the reported effect estimates and
  standard errors without a harmonised effect-size measure.
\end{itemize}

\newpage

\subsection{Askarov et al 2023
(economics)}\label{askarov-et-al-2023-economics}

\textbf{Reference:} \textcite{askarov2023significance}.

\textbf{Research question:} Impact of mandatory data sharing on
(excessive) statistical significance in economics papers

\textbf{Data availability:} The file
\texttt{Mandatory\ data-sharing\ 30\ Aug\ 2022.dta} may be downloaded
from
\url{https://github.com/anthonydouc/Datasharing/blob/master/Stata/}; the
reproducibility package is available at
\url{https://dataverse.harvard.edu/file.xhtml?fileId=7884702&version=1.0}.

\textbf{Data description and source:} The paper includes 64,000 economic
estimates surveyed by Ioannidis et al.~(2017) plus results of
``searching through numerous databases, journal websites, and also an
email survey of 109 authors known to have produced meta-analyses of
economic topics. We cast the widest net possible to capture empirical
economic estimates. Thus, some of the papers included in these
meta-analyses were originally published in political science and
psychology journals. Nevertheless, these meta-studies deal with
economics topics.'' Later they note: ``we collect 166,924 parameter
estimates from 345 distinct research areas, of which 20,121 were
published in 24 of the leading general interest and field journals.''
The paper mentions 12,521 observations from before and 2,537 from after
the introduction of data-sharing policies in 24 journals; the whole
reproducibility dataset has 32 journals (the authors removed some
smaller ones) and 22,172 rows.

\textbf{Notes:} The article by Askarov et al.~starts from the
Ioannidis-Stanley-Doucouliagos dataset, which is also used by
Arel-Bundock et al.~(another dataset in BEAR). Both datasets incorporate
additional data and target different fields of research, but there is
likely substantial overlap between the two sets.

\textbf{Data processing:} We use the same source file
\texttt{Mandatory\ data-sharing\ 30\ Aug\ 2022.dta} (11.6 MB) as in the
reproducibility package. Signed z-values were computed as reported
effect sizes divided by their standard errors.

\textbf{Study characteristics:}

\begin{itemize}
\item
  \textbf{Study ID:} we use the source study identifier.
\item
  \textbf{Meta-analysis ID:} we use \texttt{filename}.
\item
  \textbf{Topic:} we use the source's flag to distinguish macroeconomics
  or microeconomics/other.
\item
  \textbf{Method:} The source flags experimental and mixed
  observational/experimental research. We label estimates with neither
  flag set as ``not RCT, not mixed''. This category may include
  quasi-experimental research, which the source does not separately
  identify.
\item
  We do not have access to the measure of effect size (e.g.~ratio vs
  mean difference).
\end{itemize}

\newpage

\subsection{Metapsy (psychotherapy)}\label{metapsy-psychotherapy}

Studies included in \emph{Metapsy} (\url{https://www.metapsy.org/}), a
set of living meta-analytic databases of randomised trials of
psychotherapeutic interventions.

\textbf{Snapshot date:} January 2026

\textbf{Data availability:} Metapsy is maintained by an international
collaboration led by Vrije Universiteit Amsterdam. Databases are
available at \url{https://www.metapsy.org/} and through the R package
\texttt{metapsyData} \autocite{harrer2022metapsydata}. The Metapsy FAQ
states that all data provided as part of the project are open access and
that attribution is required. The source contains RCTs of
psychotherapeutic interventions.

\textbf{Data processing:} We merged 20 individual datasets in Metapsy
and performed some cleaning (e.g.~to extract study dates from titles).
Almost all estimates are Hedges' g, but for one of the databases
(\texttt{total-response}) we converted log-odds ratios to SMD (dividing
by \(\Pi/\sqrt{3}\) to scale SD of logistic regression) to keep a
comparable scale. Signed z-values were computed as effect sizes divided
by their standard errors.

\textbf{Study characteristics:}

\begin{itemize}
\item
  \textbf{Sample size and year:} \texttt{ss} is the sum of
  \texttt{n\_arm1} and \texttt{n\_arm2} when both are supplied;
  \texttt{year} is the source year or a year parsed from the study
  label.
\item
  \textbf{Study ID:} we use the source \texttt{study} label within each
  database.
\item
  \textbf{Meta-analysis ID:} we use the Metapsy database identifier.
\item
  \textbf{Topic:} we use an explicit clinical-subject lookup; datasets
  \texttt{peer-support} and \texttt{total-response} are labelled as
  ``mental health''.
\item
  Effects are standardised mean differences, mainly Hedges' g; a
  minority are converted from log odds ratios.
\end{itemize}

\newpage

\subsection{Arel-Bundock et al 2022 (political
science)}\label{arel-bundock-et-al-2022-political-science}

\textbf{Reference:} \textcite{arel2022quantitative}.

\textbf{Research question:} assess statistical power in political
science research

\textbf{Data availability:} The replication package (a zip file dated
\texttt{20241010}) is available at \url{https://osf.io/fgdet}. All data
are in the public domain.

\textbf{Data description and source:} The authors searched for
meta-analytic articles across 141 journals in political science and
contacted their authors. The resulting source dataset contains 16,649
hypothesis tests, grouped in 351 meta-analyses and reported in 46
peer-reviewed meta-analytic articles.

\textbf{Notes:} This dataset collects meta-analysis inputs, rather than
repeating data extraction from individual studies. It combines an
earlier dataset from Doucouliagos et al with another search. Therefore,
some of this dataset will overlap with Askarov et al. However, the two
articles aim to cover different fields and BEAR keeps the two datasets
separate.

\textbf{Data processing:} We used the same code as the original authors
(a replication package that uses makefiles), but skipped some additional
processing they performed for their paper. BEAR uses
\texttt{question\_id} as \texttt{metaid} (351 groups defined by
meta-analysis or research question). Signed z-values were computed as
reported estimates divided by their standard errors. We verified that
the distributions of z-values matched those in the article.

\textbf{Study characteristics:}

\begin{itemize}
\item
  \textbf{Study ID:} we use author's study ID; for \textbf{meta-analysis
  ID} we use \texttt{question\_id}.
\item
  \texttt{subset} column is the source \texttt{meta\_id}
\item
  We do not have information on the study design or measure of effect
  size.
\end{itemize}

\newpage

\subsection{Costello and Fox 2022; Yang et al 2023, 2024 (ecology and
evolution)}\label{costello-and-fox-2022-yang-et-al-2023-2024-ecology-and-evolution}

\textbf{Reference:} These two datasets come from the same replication
package and were processed using the shared pipeline in Yang et al.~2024
\autocite{yang2024large}. The original references for the datasets are
Yang et al.~2023 and Costello and Fox 2022
\autocite{yang2023publication,costello2022decline}.

\textbf{Research question:} decline effects in ecology (Costello and
Fox) and publication bias and performance of empirical research (Yang et
al)

\textbf{Data collection:} Costello and Fox examined 466 meta-analyses
identified by searching Web of Science. The reproducibility dataset
includes 88,000 rows across 466 analyses in 232 papers. Subsequently,
Yang et al.~removed duplicated studies and zeroes, discarding about 20\%
of rows. Yang et al.~collected 102 meta-analyses published in relevant
journals between 2010 and 2019, after removing records duplicated in the
Costello and Fox dataset.

\textbf{Data availability:} The two datasets may be downloaded from
\url{https://github.com/Yefeng0920/replication_EcoEvo_git/}. They are
available under a CC BY 4.0 licence.

\textbf{Data processing:} The only additional processing was to clean up
study years. Signed z-values were computed as effect sizes divided by
their standard errors. For Yang et al, BEAR uses the source file names
as meta-analysis identifiers.

\textbf{Costello and Fox:}

In \texttt{BEAR.rds}, we re-label the effect measures as SMD, response
ratio and correlation. The individual dataset keeps the original Yang et
al.~coding (Zr, lnRR, SMD).

\textbf{Yang:}

In \texttt{BEAR.rds}, we code Fisher's z as ``correlation'' and lnRR as
``response\_ratio'' for consistency with other datasets.

\textbf{Study characteristics:}

\begin{itemize}
\item
  \textbf{Study ID:} for Costello and Fox, we use \texttt{study2}; for
  Yang, we use \texttt{study\_ID}.
\item
  \textbf{Meta-analysis ID:} for Costello and Fox, we use
  \texttt{meta.analysis.id}; for Yang, we use \texttt{meta\_id}.
\item
  \textbf{Topic:} for Costello and Fox, we use ecology, evolution or
  other ecology/evolution. Yang has no topic classification.
\item
  \textbf{Costello and Fox:} the primary-study design is not classified.
\item
  \textbf{Costello and Fox:} effects are standardised mean differences,
  log response ratios, Fisher's z correlations, or less common measures.
  Less common measures include regression coefficients, log odds ratios
  and mean differences. Some standardised mean differences are recorded
  as absolute values.
\item
  \textbf{Yang:} the primary-study design is not classified.
\item
  \textbf{Yang:} effects are standardised mean differences, log response
  ratios, or Fisher's z correlations.
\end{itemize}

\newpage

\subsection{What Works Clearinghouse
(education)}\label{what-works-clearinghouse-education}

Large database of education effect sizes created by the Institute of
Education Sciences \autocite{wwc2025studyfindings}.

\textbf{Snapshot date:} 17 May 2025

\textbf{Data availability:} Flat files are available at
\url{https://ies.ed.gov/ncee/wwc/studyfindings}. The Permissions and
Disclaimers section of that website states: ``Unless stated otherwise,
all information on the U.S. Department of Education's IES website
\ldots{} is in the public domain and may be reproduced, published,
linked to, or otherwise used without permission from IES.''

\textbf{Data collection:} The WWC source data have 13,054 findings (data
rows) across 1,908 reviews. Each row is one finding, so a single source
paper can contribute multiple rows. The dataset includes effect sizes,
but not standard errors, as well as p-values reported by the studies and
p-values re-calculated by the WWC reviewers. After BEAR processing and
filtering, the main BEAR dataset includes 7,080 findings.

\textbf{Data processing:} We include only RCTs and quasi-experimental
studies (98\% of data altogether). Majority of these (about 80\%) are
RCTs.

We removed a few rows where the results were flagged as not meeting WWC
standards or as ineligible for review. We also removed rows marked by
WWC as subgroup measures to avoid treating subgroup analyses as
independent focal findings. Where available, we preferred effect sizes
and p-values calculated by WWC to the ones reported by the studies. We
replaced p-values that were equal to zero (this occurred in 3\% of the
dataset) with a truncated p-value (i.e.~assumed \(p < 10^{-16}\)). We
calculated z-values from these p-values under the assumption of
two-sided tests. We then applied the sign of the reported effect size,
so z-values are signed where the effect size is non-zero.

\textbf{Study characteristics:}

\begin{itemize}
\item
  \textbf{Study ID:} we use citation-derived author and year; we use the
  original \texttt{WWCStudy} as fall-back in some cases.
\item
  \textbf{Topic:} WWC dataset includes some useful topic
  classifications, but these are not included in \texttt{BEAR.rds} as
  they are either too granular or are overlapping flags for study
  categories
\end{itemize}

\newpage

\subsection{Cochrane Database of Systematic Reviews
(CDSR)}\label{cochrane-database-of-systematic-reviews-cdsr}

Trials with results available in CDSR,
\url{https://www.cochranelibrary.com/} \autocite{cochrane2025cdsr}.

\textbf{Snapshot date:} 20 November 2025

\textbf{Data collection:} Using the R package \texttt{cochrane}
\autocite{schwab2024cochrane}, we combined data from 8,726 CDSR review
records into a single file, \texttt{cdsr\_interventions\_19nov2025.csv}
(63 MB).

\textbf{Data processing:} We applied several processing steps and
extensive filtering.

During processing, we cleaned up study years and categorised measures
(risk ratio, odds ratio, mean difference, etc.) and experimental
designs. In particular, we created an ``RCT'' flag by scanning each
review abstract for its inclusion criteria. Most studies in CDSR are
RCTs, but some reviews also allowed quasi-experimental studies, so we
could not classify many studies. We also classify outcome rows into
\texttt{efficacy}, \texttt{safety}, \texttt{dropouts}, and \texttt{bias}
based on the comparison, outcome, and subgroup labels.

To calculate effect sizes and z-values for BEAR, we used \emph{metafor}
to calculate standardised mean differences for continuous outcomes and
probit-transformed differences for dichotomous outcomes. This puts
continuous and binary outcomes on a comparable scale.

Unlike in most of the other datasets in BEAR, we filter CDSR data
heavily. There are five important decisions:

\begin{enumerate}
\def\labelenumi{\arabic{enumi}.}
\item
  Use only the outcome and comparison from each review coded as ``1'',
  as these are most likely to be the primary outcome and most relevant
  comparison.
\item
  Keep only rows classified as \texttt{efficacy}, excluding likely
  safety, dropout, and bias-related outcomes based on comparison,
  outcome, and subgroup labels.
\item
  Use only studies with continuous and dichotomous outcomes, removing
  rows where effect estimates are based on instrumental variables or
  individual patient data.
\item
  Remove rows where the effect measure is unknown (keeping OR, RR, Peto
  OR, mean difference, standardised mean difference, and risk
  difference) and rows with no participants.
\item
  Remove count data with no events or a 100\% event rate in both arms
  (see the discussion of zero-event studies).
\end{enumerate}

These processing steps are applied in the main dataset,
\texttt{BEAR.rds}. Additional rows are still available to researchers in
the individual dataset, \texttt{Cochrane.rds}.

As a result, the distributed \texttt{BEAR.rds} file includes a smaller
analysis subset with about 36,000 rows: the first comparison and first
outcome from each review, restricted to efficacy outcomes. The processed
\texttt{data/Cochrane.rds} file is larger and is available for users who
want to make different selection decisions. It contains 760,486 result
rows, representing 90,042 review-study pairs across 6,619 Cochrane
reviews.

A large proportion of CDSR rows have binary outcomes with no events in
either arm or events in all participants in both arms. In the BEAR
analysis subset before this exclusion, 2,371 of 26,619 binary rows
(8.9\%) met one of these conditions: 2,247 rows (8.4\%) had no events in
either arm and 124 rows (0.5\%) had all participants experiencing the
event in both arms.

Zero-event studies may be important for some metascientific analyses,
but for modelling the distribution of z-values they are unhelpful. With
these rows included, the empirical distribution has extra observations
at or near zero, and the normal approximation is poor. Because of this
peak at zero, including these non-significant studies has the
counter-intuitive effect of increasing the estimated degree of
selection: the fitted distribution has a steeper drop-off between zero
and the conventional significance threshold at 1.96.

\textbf{Study characteristics:}

\begin{itemize}
\item
  \textbf{Sample size and year:} \texttt{ss} is the sum of the two arm
  totals; \texttt{year} is \texttt{study.year}.
\item
  \textbf{Study ID:} we use \texttt{study.name}.
\item
  \textbf{Meta-analysis ID:} we use the Cochrane review identifier
  (\texttt{id}).
\item
  \textbf{Topic:} we use the Cochrane specialty.
\item
  \texttt{subset} records study-data provenance: published, unpublished,
  sought or mixed.
\item
  \texttt{method}: We review meta-analysis eligibility criteria to
  determine which meta-analyses included only randomised trials. This
  classifies 69\% of estimates as RCTs; the remaining 31\% could not be
  classified reliably, although many are likely to be RCTs because
  Cochrane meta-analyses typically include RCT or quasi-randomised
  evidence.
\item
  \texttt{measure}: About 32\% of estimates are standardised mean
  differences and 68\% are differences between probit-transformed event
  proportions.
\item
  \texttt{topic} is Cochrane specialty and \texttt{subset} is source
  data type (published, unpublished, sought, or mixed);
  \texttt{outcome\_group} is also available in the processed data.
\end{itemize}

\newpage

\subsection{Chavalarias et al.~2016
(PubMed/Medline)}\label{chavalarias-et-al.-2016-pubmedmedline}

\textbf{Reference:} \textcite{chavalarias2016evolution}.

\textbf{Research question:} reporting of p-values in the biomedical
literature

\textbf{Data collection:} The authors used large-scale text mining of
MEDLINE abstracts and PubMed full texts: 4.5 million p-values in 1.6
million MEDLINE abstracts, and 3.4 million p-values in 385,000 PubMed
full-text articles.

\textbf{Data availability:} The extracted p-value dataset is publicly
available at
\url{https://dataverse.harvard.edu/dataset.xhtml?persistentId=doi:10.7910/DVN/6FMTT3}
(about 6 GB of data). According to the hosting platform, this dataset is
released under the Creative Commons Attribution-NonCommercial 4.0
International Licence.

\textbf{Data processing:} Unsigned z-values were derived from p-values
assuming two-sided tests. P-value truncation operators were processed by
collapsing variants such as \texttt{\textless{}\textless{}},
\texttt{\textless{}\textless{}\textless{}}, \texttt{\textless{}=},
\texttt{less\ than}, and \texttt{=\textless{}} into
\texttt{\textless{}}. We dropped 0.7\% of rows where p-values did not
have a ``plain'' format and 0.08\% of rows where truncation could not be
unambiguously classified. For the distributed version of BEAR 50,000
studies are selected at random, with one row per study, in order to keep
the file size manageable.

\textbf{Study characteristics:}

\begin{itemize}
\item
  \textbf{Study ID:} we use the PubMed identifier (\texttt{PMID}).
\item
  \texttt{subset} distinguishes estimates obtained by authors from
  MEDLINE abstracts vs from PMC full text.
\item
  There is no information on study design (i.e.~we cannot distinguish,
  for example, between prospective and retrospective studies or between
  RCTs and observational studies).
\item
  There is no information on effect-size measure (e.g.~SMD vs ratio).
\end{itemize}

\newpage

\subsection{SCORE (social science)}\label{score-social-science}

\textbf{Reference:} \textcite{tyner2026investigating}.

\textbf{Research question:} investigating the replicability of published
claims in the social and behavioural sciences.

\textbf{Data availability:} The SCORE replication package is available
at \url{https://osf.io/g5sny/}. The package includes replication data as
well as reviews of claims made in the included papers.

\textbf{Data description and source:} BEAR includes two SCORE outputs.
The source dataset of replications has 548 rows: 274 original claims and
274 matched replications. After BEAR processing and filtering, 267
matched replication rows are included. The source dataset for ``all
claims'' has 3,066 rows for claims from the set of replicated papers
(but no replications of these claims), where we select one statistic per
claim; after BEAR processing and filtering, 1,942 claim rows are
included.

\textbf{Data processing:} For the matched replication output, we used
the package's converted correlation scale where available for original
and replication statistics. For the set of all claims, effects are
heterogeneous across the source papers.

When several inputs are available for constructing a z-value in the set
of all claims, we prefer reported z, then reported t, then coefficient
divided by standard error, then signed square-root F for numerator df 1,
then a 95\% confidence interval with a point estimate, then a two-sided
p-value conversion. When the selected statistic only supplies a p-value
and no sign can be inferred, the z-value is unsigned. In the text for
all claims, it is typical to have multiple statistics backing up a
single claim, e.g.~``F(1,38) = 3.73, p = .033, partial eta-squared =
.16; F(1,39) = 7.28, p = .010, partial eta-squared = .16; F\textless1''.
We created a rule to pick one statistic per claim using agreement with
the SCORE classification as significant or non-significant, the
statistic's provenance, exactness, the available effect and standard
error, and text order.

\textbf{Study characteristics:}

\begin{itemize}
\item
  \textbf{Study ID:} we use the DOI/source-paper identifier in both
  outputs.
\item
  \textbf{Topic:} field of study (e.g.~``political science'') as
  classified by the authors.
\item
  The underlying study design is not systematically described.
\item
  Measures: SCORE ``claims'' dataset measures include regression
  coefficients, standardised mean differences, correlations, and
  variance-explained statistics; but the effect measure is unclassified
  for over 80\% of data. SCORE replications use correlations where
  available.
\end{itemize}

\newpage

\subsection{Lang 2025 (economics)}\label{lang-2025-economics}

\textbf{Reference:} \textcite{lang2025credible}.

\textbf{Research question:} How frequently are statistically significant
rejections of null hypotheses in recent empirical economics false
positives?

\textbf{Data collection:} Lang samples empirical articles from the top
five economics journals in 2021 and also articles previously collected
in Brodeur, Cook, and Heyes, who sampled articles published in 2015 and
2018 in 25 leading economics journals. (See also the ``Brodeur'' dataset
in BEAR for more details.) The sample consists mainly of natural,
laboratory, or field experiments, excluding articles focused on theory,
structural estimation, or econometric methods.

Lang re-extracted results from all articles rather than using the
hypothesis tests extracted by Brodeur, Cook, and Heyes. Research
assistants read the abstracts and results text and retained estimates
emphasised in the abstract or described by the authors as principal
results. When multiple specifications tested the same hypothesis, they
retained the authors' preferred specification or the specification whose
magnitude was used in subsequent analysis. Robustness checks and
heterogeneity analyses were excluded unless the heterogeneity analysis
was itself central to the article.

The source data contain 3,885 extracted hypothesis tests from 736
articles. Lang's main sample in the cited article restricts these to
rejected null hypotheses with \(1.96 \leq |t| < 10\), leaving 2,082
hypothesis tests from 663 articles. The dataset therefore does not
identify one focal hypothesis per article. An article can contribute
several hypothesis tests. However, only one preferred specification is
retained for each hypothesis. (Lang paper also reports analyses
restricted to one null hypothesis per article. The supplementary
material contains a separate vector with one t-statistic for each of the
736 articles, but this vector does not contain row or article
identifiers that allow it to be linked directly to the 3,885-test table.
We therefore do not use this in BEAR data.)

\textbf{Data processing:} We use the 3,885 results. We use Lang's
supplied, derounded t-statistic, derounded coefficient
(\texttt{coef\_dr} in the source data) and standard error
(\texttt{se\_dr}); the latter two are available in about 95\% of cases.
(Using either coefficients or t-statistics to calculate z-values leads
to practically identical results; we use the latter.) We calculate the
corresponding two-sided normal-approximation p-value. We retain some
additional columns in \texttt{Lang.rds}.

We also pull DOIs for source articles using bibliographic metadata
matching. These DOIs were obtained by us from Crossref, rather than
supplied by Lang; the saved assignments are preserved when reprocessing
the dataset. This is done automatically via Crossref. The matches are
very good in almost all cases. For several edge cases we used LLM review
to adjudicate.

\textbf{Study characteristics:}

\begin{itemize}
\item
  \textbf{Study ID:} \texttt{studyid} is \texttt{Lang\_paper\_} followed
  by Lang's \texttt{unique\_paperid}, but with a few minor corrections
  from us (see GitHub repository for details). Article DOIs are also
  available as an additional column in \texttt{Lang.rds}.
\item
  \textbf{Subset:} \texttt{subset} distinguishes observations drawn from
  the Brodeur--Cook--Heyes 2015/2018 source universe from Lang's 2021
  top-five additions.
\item
  \textbf{Method:} The \texttt{method} labels supplied with the Lang
  data mix study designs, identification strategies, and estimators.
  Large majority of estimates are one of ``RCT'', ``DID'', ``IV'',
  ``RD'', and ``OLS''. \texttt{Lang.rds} retains the original method
  labels. When constructing \texttt{BEAR.rds}, we collapse some label,
  e.g.~\texttt{RCT\ (DID)} and \texttt{RCT\ (IV)} become \texttt{RCT};
  \texttt{DID\ (matching)} becomes \texttt{DID}; and \texttt{RD\ (DID)}
  and \texttt{RD\ (IV)} become \texttt{RD}. For estimates labelled
  \texttt{OLS}, \texttt{matching}, \texttt{SYSTEM-GMM}, or
  \texttt{IV-DID\ (matching)}, we retain the estimates but set
  \texttt{method} to missing.
\item
  The dataset does not have data on type of measure to allow us to
  classify estimates as in some other datasets in BEAR (e.g.~risk risk
  vs standardised mean difference).
\end{itemize}

\newpage

\subsection{Szucs and Ioannidis 2017 (cognitive
neuroscience)}\label{szucs-and-ioannidis-2017-cognitive-neuroscience}

\textbf{Reference:} \textcite{szucs2017empirical}.

\textbf{Research question:} empirical assessment of published effect
sizes, statistical significance, and power in the literature on
cognitive neuroscience and psychology.

\textbf{Data availability:} The article and supporting data are
available from PLOS Biology at
\url{https://journals.plos.org/plosbiology/article?id=10.1371/journal.pbio.2000797}.

\textbf{Data description and source:} The source contains 26,841
text-mined t-test records from 3,801 papers in 18 journals (cognitive
neuroscience, psychology, and medical). The processed dataset preserves
all three fields; in the main BEAR dataset, we include only the
cognitive neuroscience subset (since the psychology and medical
literature are well covered by other datasets in BEAR).

\textbf{Notes:} The extraction targeted t-test records reported in
article texts and did not mine tables.

\textbf{Data processing:} We followed our default procedure for deriving
z-values from t-values.

\textbf{Study characteristics:}

\begin{itemize}
\item
  \textbf{Study ID:} to make look-ups of articles easier, we use
  \texttt{journal\_\textless{}journal\_code\textgreater{}\_article\_\textless{}article\_id\textgreater{}}.
\item
  \textbf{Topic:} we use the source \texttt{field} classification; note
  that \texttt{BEAR.rds} contains only cognitive neuroscience rows.
\item
  The source data does not classify study designs.
\item
  The original dataset labels all results as coming from a t-test, but
  there is no additional information on what was the underlying measure.
\end{itemize}

\newpage

\subsection{ClinicalTrials.gov
snapshot}\label{clinicaltrials.gov-snapshot}

Trials that reported results at ClinicalTrials.gov
\autocite{clinicaltrials2025}.

\textbf{Snapshot date:} 12 August 2025

\textbf{Data source:} This is a set of flat files of studies obtained
from the AACT downloads page
(\url{https://aact.ctti-clinicaltrials.org/downloads/}; over 2 GB of
data for relevant tables), with relevant columns extracted and joined by
us before BEAR processing.

\textbf{Data availability:} The terms and conditions for
ClinicalTrials.gov are at
\url{https://clinicaltrials.gov/about-site/terms-conditions}. We make
distributions of derived effect sizes available in accordance with the
terms and conditions. We only use publicly available data and do not
imply any endorsement of this work by the National Library of Medicine,
ClinicalTrials.gov or AACT.

\textbf{Data processing:} Processing was extensive. The full data table
(\texttt{BEAR\_data/clinicaltrialsgov.rds}) differs in some respects
from the main BEAR dataset. We briefly discuss these differences here.
Because of its rich structure, please refer to the data dictionary for
the full data table.

Briefly, BEAR targets primary outcome results for completed studies
using two data sources:

\begin{enumerate}
\def\labelenumi{(\arabic{enumi})}
\item
  Author-reported effect sizes
\item
  Effect sizes derived from raw outcome data (group means, counts),
  alongside source data
\end{enumerate}

Where both (1) and (2) are available, we prefer (1), but also make (2)
available in the full data table. For (2), we use standardised mean
differences for continuous outcomes and probit differences for count
outcomes.

When working with the full dataset, variable
\texttt{author\_raw\_overlap\_key} can be used to link the two different
calculations of the same effect; variable \texttt{include\_in\_bear}
identifies rows included after applying this rule preferring
author-reported estimates.

The main BEAR dataset keeps studies that (1) are randomised, (2) are
interventional, and (3) have fewer than 20 candidate effect rows (see
the description in this section).

The full ClinicalTrials.gov table uses only completed studies but does
not apply these three restrictions. It contains 151,313 rows from 27,446
studies; 46.3\% are author-reported and 53.7\% are derived from raw
data. The randomised interventional subset used in the main BEAR release
contains 60,470 rows from 23,060 studies.

Note that a single ClinicalTrials.gov study can contribute several
primary outcome IDs, and each primary outcome can contain many
measurements. The full data table records this as
\texttt{n\_effect\_rows\_per\_study} and
\texttt{n\_outcomes\_per\_study}. The median study contributes one
outcome ID and two rows to the full data table. Studies with 20 or more
candidate effect rows remain in the full data table but are excluded
from the main BEAR build.

In the full data table, some rows derived from raw data have
\texttt{direction\_unknown\ ==\ TRUE}, meaning that the
control/reference arm could not be reliably inferred, so the sign should
be treated as random. Users interested in controlled trials or signed
z-values should filter the results accordingly.

See \emph{Standard procedure for dealing with p-values and confidence
intervals} for how we created and chose the z-values based on p-values
and confidence intervals. Briefly, we calculated z-values using two
separate procedures (based on p-values and confidence intervals) and, if
both were available, we made a judgement on which one was more
appropriate using a predefined rule. ClinicalTrials.gov reports the
number of sides for confidence intervals but not for p-values, so we
treated p-values as two-sided. We dropped rows only when \texttt{z} was
missing. Notably, infinite z-values can arise when \texttt{p\_value} is
recorded as exactly 0, and these are included because they are not
missing.

\textbf{Study characteristics:}

\begin{itemize}
\item
  \textbf{Study ID:} we use the ClinicalTrials.gov trial identifier
  (\texttt{nct\_id}).
\item
  \texttt{subset} records trial phase.
\item
  We include only randomised interventional studies, including parallel,
  crossover and factorial designs.
\item
  About half of all effects (by which we mean both reported and our
  derived estimates) are standardised mean differences, mean
  differences, or probit differences. Another 14\% are hazard, odds,
  risk, geometric, or rate ratios; 5\% have the other category (in
  \texttt{BEAR.rds} this is coded as ``other''), and 20\% have no
  measure.
\end{itemize}

\newpage

\subsection{EU Clinical Trials Register (EU
CTR)}\label{eu-clinical-trials-register-eu-ctr}

Trials that reported results in the EU Clinical Trials Register,
\url{https://www.clinicaltrialsregister.eu/} \autocite{euctr2025}.

\textbf{Snapshot date:} 18 January 2026

\textbf{Data collection:} We used the R package \texttt{ctrdata}
\autocite{herold2025ctrdata} to create a snapshot of EU CTR
(\texttt{collection\ =\ "euctr"}); we downloaded only trial records with
results. We stored the downloaded records in a local SQLite database and
extracted a fixed set of protocol and results fields plus several fields
computed by \texttt{ctrdata} (trial phase, sample size, and a ``first
primary endpoint'' p-value/size derived by the package). We use only the
EU CTR portion of this extraction in BEAR.

\textbf{Data availability:} EUCTR is publicly accessible. The EUCTR
legal notice notes that publication on EUCTR does not constitute an
endorsement of the information reported.

\textbf{Data processing:} Processing was extensive. We use only primary
endpoints. We default missing CI levels to 95\% and where truncation of
p-values was not stated we assume equality (the latter is the case for
about 15\% of data). Where the data do not report whether a test is one-
or two-sided, we assume two-sided. The rest of the derivation of
z-values follows the \emph{Standard procedure for dealing with p-values
and confidence intervals}. We set \texttt{year} to the year of entry
since completion dates were unavailable. We dropped rows only when
\texttt{z} was missing; infinite z-values are included. We harmonise
phase labels to follow the same naming convention as ClinicalTrials.gov.

\textbf{Study characteristics:}

\begin{itemize}
\item
  \textbf{Sample size:} \texttt{ss} is the extracted endpoint sample
  size (\texttt{n}).
\item
  \textbf{Study ID:} we use the EudraCT trial identifier (\texttt{id}).
\item
  \texttt{subset} records trial phase.
\item
  The registry extract does not provide a row-level randomisation
  indicator, so BEAR does not classify the trial design.
\item
  Effect measures are derived from reported estimand labels; about half
  of these are mean differences, hazard ratios, odds ratios, risk
  differences, or risk ratios. Another one third have other category (in
  \texttt{BEAR.rds} this is coded as ``other''), and 23\% have no
  measure.
\item
  \texttt{subset} records registry phase.
\end{itemize}

\newpage

\subsection{Barnett and Wren 2019
(PubMed/Medline)}\label{barnett-and-wren-2019-pubmedmedline}

\textbf{Reference:} \textcite{barnett2019examination}.

\textbf{Research question:} Bias for statistical significance in health
and medical journals

\textbf{Data availability:} The dataset \texttt{Georgescu.Wren.RData}
may be downloaded from
\url{https://github.com/agbarnett/intervals/tree/master/data} or
\url{https://github.com/jdwren/ASEC}. The repositories give no clear
licence, but the BMJ Open article is CC BY 4.0.

\textbf{Data description and source:} The dataset is a collection of
confidence intervals for ratio estimates from Medline papers published
between 1976 and 2019. Authors scraped (via regular expressions,
followed by an independent check using another data mining algorithm,
manual checks for 10,000 intervals) 968,000 CIs from abstracts and
350,000 from full texts.

\textbf{Notes:} These are ``ratio estimates'' such as odds ratios,
hazard ratios and risk ratios. Note that binary outcomes tend to contain
(much) less information than continuous outcomes.

\textbf{Data processing:} We exclude a small fraction of data. First, we
restrict the data to 95\% confidence intervals: where \texttt{ci.level}
was missing we assume it was 0.95; among intervals with a known
\texttt{ci.level}, about 0.3\% were not 95\% and are dropped. We drop
intervals with non-positive width
(i.e.~\texttt{lower\ \textgreater{}=\ upper}). We use the log scale (we
replace zero or negative lower bounds with a small positive constant;
this affects about 0.1\% of the sample) and back out an approximate
standard error and point estimate under the usual normal approximation
for a 95\% CI, setting the point estimate to the midpoint of the
interval.

For the main version of BEAR dataset, \texttt{BEAR.rds}, 50,000 studies
are selected at random, with one row per study, in order to keep the
file size manageable.

\textbf{Study characteristics:}

\begin{itemize}
\item
  \textbf{Study ID:} we use the PubMed identifier (\texttt{pubmed}).
\item
  There is no information on study design (i.e.~we cannot distinguish,
  for example, between prospective and retrospective studies or between
  RCTs and observational studies).
\item
  Effects are ratios on a logarithmic scale, without distinguishing
  odds, risks, hazards etc.
\end{itemize}

\newpage

\subsection{Open Science Collaboration 2015
(psychology)}\label{open-science-collaboration-2015-psychology}

\textbf{Reference:}
\textcite{opensciencecollaborationEstimatingReproducibilityPsychological2015a}.

\textbf{Research question:} replications of a quasi-random sample of 100
experiments in psychology.

\textbf{Data collection:} The authors make a comprehensive dataset of
results available with their paper. BEAR includes only the replications
because all original studies in that set have \(p < 0.05\).

\textbf{Data availability:} The data file \texttt{rpp\_data.csv} is
available in the paper's repository,
\url{https://osf.io/ytpuq/overview}, under a CC0 1.0 Universal licence.

\textbf{Data processing:} We derived z-values from the replication
p-values assuming two-sided p-values (we use the p-value column for the
replication, not the p-value for the original study). We applied the
sign of the replication correlation, so z-values are signed where that
correlation is available. We treated p-values recorded as exactly 0 as
truncated and recorded this as a lower bound on z (z-operator
\texttt{\textgreater{}}); otherwise we treated p-values as exact
(z-operator \texttt{=}). We retained the replication sample size as
\texttt{ss}.

In \texttt{BEAR.rds}, the measure is recorded as ``correlation'', based
on the fact that is the column from original dataset that we use.

\textbf{Study characteristics:}

\begin{itemize}
\item
  BEAR treats the replication studies as RCTs.
\item
  The retained effect measure is the replication correlation.
\end{itemize}

\newpage

\section{Standard procedure for dealing with p-values and confidence
intervals}\label{standard-procedure-for-dealing-with-p-values-and-confidence-intervals}

Most datasets in BEAR come with pre-computed effect sizes and standard
errors. Some, especially databases of clinical trials and large sets of
scraped studies, report p-values. Moreover, reporting in clinical trial
databases is done by investigators, and the data-entry fields are not
standardised, requiring more extensive clean-up and a procedure for
deriving z-values. Lastly, in some cases datasets contain confidence
intervals and effect sizes, but z- and p-values may be missing.

\subsection{Data cleaning and default
values}\label{data-cleaning-and-default-values}

For confidence interval levels, we set implausible values to missing
(e.g.~values above 100 or below 0) and standardise proportions and
percentages: investigators sometimes report 90 and sometimes 0.9 for a
90\% confidence interval. In EU CTR, confidence levels can be missing,
and we assume 95\% by default. Some data state whether CIs or p-values
are one- or two-sided.

Similarly, when working with p-values we assume a two-sided test when no
additional information is stated.

\subsection{Confidence intervals}\label{confidence-intervals}

When working with CIs, we analyse ratio measures on the log scale when
the lower bound, upper bound, and point estimate are all strictly
positive; we keep non-ratio measures on the raw scale.

To derive SEs from confidence intervals, we first compute the critical
value. To avoid undefined critical values (e.g.~at 100\% or extremely
close to it), we bound the implied \(\alpha\) at a small value. We then
compute two one-sided standard error estimates from the upper and lower
bounds and use their average when both are finite.

We compute a simple symmetry diagnostic for the confidence interval on
the chosen scale, defined as the smaller of the two half-widths divided
by the larger. We treat intervals as ``symmetric enough'' only if this
ratio exceeds 0.8. This arbitrary threshold affects when we use the
CI-derived z-statistic; otherwise, we use a p-derived z-statistic where
one is available.

Where additional information is available, we look for potentially
non-Wald intervals by matching text in the parameter-type column. We
search for keywords such as median, Hodges, posterior/Bayes, exact,
Fieller, bootstrap, and permutation. For flagged estimates, we prefer to
derive the z-value from the p-value.

\subsection{p-values}\label{p-values}

When a source reports t-statistics and degrees of freedom, we first
derive the two-sided p-value from the t distribution and then derive the
z-value (retaining the sign of the original t-statistic, if available).

For clinical trials databases (ClinicalTrials.gov and EU CTR), we
sometimes find ambiguity in how p-values are recorded. Some significant
results report p-values close to 1, suggesting that authors report the
cumulative distribution, \(\Phi(z)\). Therefore, for each row we also
calculate \(2\cdot\min\{p,1-p\}\). When a CI-derived z-value is
available, we choose the p-value that gives the closest z-value.

To avoid dropping rows, we retain unsigned z-statistics when the sign
cannot be determined. If the sign of the effect is known, we apply it to
the p-derived z; otherwise, we use its absolute value.

\subsection{Choice of a single z-value when multiple candidates are
available}\label{choice-of-a-single-z-value-when-multiple-candidates-are-available}

In situations where \(z\) can be derived from both CIs and p-values, we
choose a single ``best'' z-statistic per row as follows. We prefer the
CI-derived \(z\) when it looks like a plausible Wald interval on the
chosen scale, meaning: (1) finite \(z\) and standard error, (2) positive
standard error, (3) no non-Wald flag from a keyword search of the
measure label, and (4) either missing symmetry information or symmetry
above 0.8. Otherwise, we use the p-derived z.

\subsection{Truncation and infinite
z-values}\label{truncation-and-infinite-z-values}

For reported two-sided p-value bounds, \(p<c\) gives a lower bound on
\(|z|\). BEAR records this inequality using \texttt{z\_operator}. The
operator applies to the magnitude, including when the stored z-value is
negative. For example, a stored value of \(-1.96\) with operator
\texttt{\textgreater{}} means a negative direction and magnitude greater
than 1.96, not \(z>-1.96\).

Exact zero p-values can produce infinite z-values. We retain infinities
in BEAR source data, although signal-to-noise models truncate these
infinite values.


\end{document}
